\section*{Overview}

KMP is the classical linear-time pattern-matching algorithm built on the prefix function. I like it because the core
idea is small: when a mismatch happens, do not restart from scratch. Reuse the longest border of what has already
matched.

\section*{When to Use It}

Use it when:

\begin{itemize}
  \item you need all occurrences of a pattern in a text,
  \item the alphabet can be arbitrary and hashing is unnecessary or too fragile,
  \item you want a deterministic linear-time matcher,
  \item the prefix function itself is useful for border or periodicity questions.
\end{itemize}

For single ad hoc substring checks, the standard library may be enough. KMP becomes interesting when the linear
guarantee or the prefix function structure matters.

\section*{Core Idea}

For each prefix ending at position \(i\), the prefix function \(\pi[i]\) stores the length of the longest proper
prefix that is also a suffix of the substring ending at \(i\).

When matching a pattern against a text, suppose the first \(k\) pattern characters already match and the next
character mismatches. Instead of throwing away all progress, KMP jumps to the longest border of the matched prefix,
whose length is \(\pi[k - 1]\). That is exactly the longest candidate that could still extend to a full match.

\section*{Operations}

\begin{itemize}
  \item compute the prefix function of the pattern,
  \item scan the text left to right,
  \item on mismatch, repeatedly jump with \(\pi\),
  \item on full match, report the occurrence and continue from the next border.
\end{itemize}

The same prefix-function routine is also useful by itself for border arrays, repetitions, and string compression
questions.

\section*{Correctness Intuition}

At every point in the scan, KMP keeps the invariant that \(\texttt{matched}\) is the length of the longest prefix of
the pattern that matches a suffix of the text processed so far.

When a mismatch occurs, any longer candidate is impossible, so the only sensible fallback is the longest border of
that matched prefix. Repeating this step eventually finds the longest viable suffix-prefix overlap. No text character
needs to be revisited from scratch, which is why the total work stays linear.

\section*{Complexity Analysis}

\begin{itemize}
  \item prefix function: \(O(n)\) for pattern length \(n\),
  \item text scan: \(O(m)\) for text length \(m\),
  \item total matching time: \(O(n + m)\),
  \item memory: \(O(n)\).
\end{itemize}

\section*{Implementation}

The note code exposes two small functions:

\begin{itemize}
  \item \texttt{prefix\_function(pattern)}
  \item \texttt{kmp\_search(text, pattern)}
\end{itemize}

That is usually the minimal contest-ready surface area. If I only need border information, I stop after the first
function.

\section*{Common Pitfalls}

\begin{itemize}
  \item Confusing the prefix function with the Z-function.
  \item Forgetting that \(\pi[i]\) is a \emph{proper} prefix length.
  \item Mishandling the empty pattern edge case.
  \item Writing the fallback loop incorrectly and skipping valid borders.
  \item Reaching for rolling hash when deterministic linear matching would be simpler.
\end{itemize}

\section*{Practice Problems}

\begin{itemize}
  \item Find all pattern occurrences in a long text.
  \item Count borders or detect periodic strings with the prefix function.
  \item Problems that combine pattern matching with automaton-style DP over prefixes.
\end{itemize}

\section*{References}

\begin{itemize}
  \item \href{https://cp-algorithms.com/string/prefix-function.html}{cp-algorithms: Prefix function and KMP}.
  \item \href{https://atcoder.github.io/ac-library/production/document_en/string.html}{AtCoder Library string documentation}.
  \item \href{https://usaco.guide/CPH.pdf}{Competitive Programmer's Handbook}.
\end{itemize}
